\section{Simulator}
\label{sec:simulator}

We describe a simulator $\mathcal S$ that runs the real-world adversary
while interacting with $\FMPC$.  It simulates the public ledger, the
universal-sampler parameters and oracle, and the information disclosed
on corruption.  It knows each public round circuit, the wire assignments,
and the communication schedule.  It does not receive honest parties'
input or output bits, nor the functionality's secret state.  An allowed
corruption supplies the affected party's retained inputs and already-computed
outputs, including outputs pending private delivery.
Honest inputs arrive at the public round start, so their arrival times
need not be inferred from their values.

The main separation is between preparing a garbling and releasing its
input encoding.  During preparation, $\mathcal S$ produces a simulated
garbling and non-committing ciphertexts.  Later it chooses the public
masked outputs and an encryption of a zero state.  These public values
suffice to obtain one simulated input encoding.  Share completion and
receiver explanations then make the selected committees release exactly
that encoding.  This section specifies these actions; the
indistinguishability argument is separate.
Whenever a nonce or label committee is assigned, $\mathcal S$ checks
its number of initially corrupt members and aborts if this exceeds the
sharing threshold $t$.  Below we describe the execution when these
checks pass.

\begin{description}
  \item[Registration and retained state.]
    Simulate the initial random corruption experiment for computation
    parties.  Generate genuine keys for initially corrupt parties and
    disclose their prescribed state.  For each initially honest delivery
    key, run
    \[
      (\widetilde{\pk},\widetilde c,\SimState_P)
        \sample\SimNCE(1^\lambda).
    \]
    Register $\widetilde{\pk}$, but retain $\widetilde c$ and
    $\SimState_P$ privately until the key is assigned its unique delivery
    role.  Thus registration can precede that assignment by arbitrarily
    many rounds.  Generate input-authorization strings $S_0,S_1$, their
    public images, and their random tapes honestly.

    To explain a delivery of a message $m$ to $P$, compute
    $\widetilde\rho\sample\ExplainNCE(\SimState_P,m)$ and obtain its
    secret key by rerunning $\KeyGen(1^\lambda;\widetilde\rho)$.
    Save this explanation once made: every subsequent disclosure uses
    the same key and complete key-generation tape.
    If the key is assigned to a nonce committee, make this explanation
    for a fixed dummy message of the prescribed length, discard the
    unpublished special ciphertext, and use ordinary encryptions for
    all nonce-share packets under that key.
    If an output recipient is corrupted before its pad ciphertext has
    been included in a sample, likewise explain its key using a dummy
    message and discard the unpublished special ciphertext.  Disclose
    the resulting key, tape, and retained history; its later pad
    ciphertext will be an ordinary encryption under this fixed key.

  \item[Public sampling.]
    Maintain a table indexed by exact distribution descriptions.
    A service $\SimSample(d)$ returns the stored sample if one exists.
    Otherwise, for a designated main-round description, it prepares the
    simulated sample described below.  For any other description,
    including the nonce-only bootstrap, it samples fresh uniform coins
    $r_d$ and computes $d(r_d)$ honestly.  It stores the result before
    returning it.

    Use this service as the oracle for the universal sampler's simulator:
    \[
      (U,\SimState_U)
        \sample\SimSetup^{\SimSample}(1^\lambda).
    \]
    Put $U$ in the simulated setup.  On a public oracle query $q$, answer
    with $a$, updating the retained state according to
    \[
      (a,\SimState_U)
        \sample\SimRO^{\SimSample}(\SimState_U,q).
    \]
    The adversary still runs the public $\Sample$ algorithm locally;
    it is not restricted to receiving samples from a black-box interface.
    Here $\SimSample$ specifies the source supplied to $\SimSetup$ and
    $\SimRO$, not a new programming operation assumed of the sampler.

  \item[Identify the next description before nonce release.]
    Once the protocol freezes the next round's circuit, assignments,
    delivery keys, and committee rosters, designate its exact description
    $d_s$, including its already-fixed nonce $Q_s$, in the sampling service.
    The simulator knows this nonce before its public release.  For the
    bootstrap it obtains $Q_1$ by decrypting enough nonce-share packets
    with its retained nonce-committee keys and reconstructing; this does
    not require the ideal sample source to reveal a sampling tape.
    Subsequent nonces are generated by $\mathcal S$ in the preceding
    main-round sample.

    If $d_s$ was already answered as an ordinary description, abort;
    never replace an earlier table entry.  Otherwise reserve that entry
    for the simulated main-round sample.  Simulate the nonce committee's
    scheduled key disclosures and subsequent corruptions using its
    already-fixed keys and tapes.  Robust reconstruction releases the
    one pre-existing $Q_s$, even if corrupt members withhold or alter
    their contributions.  The other fields remain fixed after release.

  \item[Prepare a main-round sample.]
    Generate a genuine state-transfer pair
    $(\statepk_s,\statesk_s)$.  Let $C_s$ denote the augmented round
    circuit: it receives the preceding state ciphertext and the
    successor's public key as public inputs, performs the application
    computation, and outputs masked application bits and an encrypted
    new state.  The first circuit starts from the prescribed empty state
    rather than decrypting a predecessor's ciphertext.  Use its topology
    to compute
    \[
      (\GC_s,\SimState_s)
        \sample\SimGarble(1^\lambda,\Topo(C_s)).
    \]
    Include $\statepk_s$ and $\GC_s$ in the public sample.
    The circuit layout is fixed independently of the hardcoded secret
    key, output pads, and encryption coins.  Constructing its topology
    therefore does not require assigning those secret values.

    Include the retained special ciphertext for each still-honest
    input-provider, output-recipient, and label-committee delivery role.
    For an already-corrupt output recipient, choose a uniform pad $r_i$,
    include its genuine encryption under the recipient's registered key,
    and retain $r_i$ for the public-output step.
    For each label committee, fix uniform shares at its initially corrupt
    positions and genuinely encrypt their packets.  Fix the public
    activation instructions and default-committee flags as in the protocol.
    Independently sample $Q_{s+1}$, robustly share it, and genuinely encrypt
    its shares to the reserved successor nonce committee.  These packets
    are part of the sample and are available before evaluating $\GC_s$.
    Retain the garbling simulator state for the later input-encoding call.

  \item[Simulate input disclosures.]
    For each honest input provider choose an independent uniform
    committee index $j\in\bits$ and post $S_j$ in the input phase.
    This fixes its selected committee without requiring its input bit.
    If the provider is subsequently corrupted, obtain its accepted bit
    $b$ from $\FMPC$ and explain its earlier delivery ciphertext as
    containing
    \[
      \pi_w=b\oplus j.
    \]
    The explained private computation then posts
    $S_{b\oplus\pi_w}=S_j$, as already recorded.  Supply the saved
    authorization randomness and the explained receiver key, tape, and
    private-call history.  Do not submit a replacement input to $\FMPC$:
    the honest input was already accepted, and later corruption cannot
    change it.

    Under the present corruption rule a fresh honest input provider
    cannot be corrupted before this post.  Consequently the simulator
    does not choose its input on the adversary's behalf.  Initially corrupt
    computation parties may send arbitrary messages; process those
    messages and the common deadlines exactly as the protocol prescribes.

  \item[Fix the public evaluation result.]
    After input choices close, the ideal round has computed its outputs,
    and the successor sample is available, choose an independent uniform
    masked bit $z_i$ for each still-honest output recipient.
    No output bit $y_i$ is needed for this choice.
    For an already-corrupt recipient, obtain its output $y_i$ from
    $\FMPC$ and set $z_i=y_i\oplus r_i$, using its previously fixed pad.
    The outputs of honest recipients remain unknown to $\mathcal S$.
    Choose fresh encryption coins and form
    \[
      c_s\sample\Enc(\statepk_{s+1},0^{|\sigma_s|}),
      \qquad
      X_s\sample\SimEncode(\SimState_s,(z,c_s)).
    \]
    Here $z$ is the vector of masked bits and $|\sigma_s|$ is the
    publicly specified length of the updated state.
    Make this one encoding call before any honest label committee
    speaks, and keep its result fixed.  The same public $c_s$ is the
    incoming state ciphertext for round $s+1$; simulating that round
    does not require decrypting it.

  \item[Release the selected labels.]
    Parse $X_s$ into its selected wire labels and common material
    $\GarbleAux$.  For each selected committee, use $\CompleteShares$
    to complete its previously fixed corrupt shares to the corresponding
    label bundled with $\GarbleAux$.  This applies both to private-input
    committees selected by preimages and to public-input committees
    selected by the actual incoming ciphertext and successor public key.

    For each honest member, explain its delivery ciphertext as the
    resulting share and activation instructions \emph{before} posting
    the explained decryption key.  On a later corruption, reproduce this
    same key, tape, decrypted packet, and retained computation history.
    The unselected committees' honest members remain silent and protected;
    their initially corrupt members retain the shares already fixed.
    Simulate reconstruction and public evaluation with the released
    encoding.

  \item[Private outputs and further rounds.]
    Private application outputs are delivered by $\FMPC$, not computed
    by $\mathcal S$.  Handle chosen output-recipient corruptions whenever
    they occur.  If a recipient is first corrupted after its pad
    ciphertext is included in a sample but before $z_i$ is fixed,
    choose a uniform pad $r_i$ and explain the ciphertext accordingly.
    Later use its ideal output to set $z_i=y_i\oplus r_i$.
    If the recipient is first corrupted after $z_i$ has been fixed,
    obtain its output $y_i$ from $\FMPC$ and explain the pad ciphertext as
    \[
      r_i=z_i\oplus y_i.
    \]
    The ideal output is already determined at this point by the ordering
    above.  In every case, disclose the explained receiver key and complete
    tape, together with the recipient's retained private computations.
    An explanation is never redrawn, and an already-fixed $z_i$ is never
    changed.  Corruption before the pad ciphertext is sampled is handled
    by the registration rule.

    Continue with the already-prepared successor, preparing one further
    round before its evaluation.  Internal identity-buffer rounds have
    no application inputs or outputs: simulate their state-ciphertext
    transfer without making an additional invocation of $\FMPC$.
    Keep all explanations for parties whose states have been exposed,
    whether by posting or corruption, but do not
    disclose a garbler tape or state-transfer secret key as a party's
    state: no party holds that material in the real protocol.
\end{description}
